%% 6-9 複素平面：かけ算は回転と拡大（第14回：組 V6-5。ページ制限なし。偶数ページ）
%% ほかの巻が指している形：6-8 の「j をかけると90°回る」を、かけ算一般（大きさはかけ算、角はたし算）へ広げる。
%%   4-9 の加法定理（道すじごと回す）＝複素数のかけ算、4-15 の極座標と tan^{-1} の区画、4-7 の2リンクのうで（同じ数 40・30・30°・60°）。
%%   偏角の範囲は −180°＜θ≦180°（Python の cmath.phase、表計算の IMARGUMENT・ATAN2 と同じ）。4-15 の 0°〜360° とは一般角（4-8）で同じ向き。
%%   1の n 乗根は もっと知りたい（6-8 例題8(3) の s^2+s+1=0 の解、三相交流の120°）。r∠θ は専門書。
\makeatletter
\newcommand\bkvijhead[1]{\par\addvspace{6pt}\Needspace{6\baselineskip}\noindent{\sffamily\bfseries\color{\bkTT}#1}\par\nobreak\vspace{2pt}}
\newcommand\bkvijbox[1]{\par\addvspace{3pt}\noindent\textbf{#1}\par\noindent\ignorespaces}
\newcommand\bkvijframe[2][100mm]{\par\addvspace{4pt}{\centering\tikz\node[draw=\bkTT, line width=0.9pt, fill=bkTheme!10, rounded corners=1.5mm, inner xsep=6mm, inner ysep=2.5mm, align=center]{#2};\par}\addvspace{4pt}\noindent\ignorespaces}
\newcommand\bkvijline[1]{\par\addvspace{3pt}{\centering #1\par}\addvspace{3pt}\noindent\ignorespaces}
\newcommand\bkvijq{\par\vspace{2pt}\noindent\hangindent=4.3\zw\hangafter=1\hspace*{2.6\zw}}
\newcommand\bkvija{\par\vspace{2pt}\noindent\hangindent=2.6\zw\hangafter=0\relax}
\newcommand\bkvijb{\par\vspace{2pt}\noindent\hangindent=4.3\zw\hangafter=0\relax}
% 複素平面の軸（実軸・虚軸の名前つき）  \bkvijaxes{左}{右}{下}{上}
\newcommand\bkvijaxes[4]{%
  \draw[bkGray, line width=0.6pt, -{Stealth[length=4pt]}] (#1,0) -- (#2,0) node[anchor=north east, font=\sffamily\footnotesize, text=bkInk, inner sep=1pt, yshift=-1mm] {実軸};
  \draw[bkGray, line width=0.6pt, -{Stealth[length=4pt]}] (0,#3) -- (0,#4) node[anchor=north west, font=\sffamily\footnotesize, text=bkInk, inner sep=1pt, xshift=1mm] {虚軸};}
% 原点からの矢印  \bkvijarrow{色}{(x,y)}
\newcommand\bkvijarrow[2]{\draw[#1, line width=1.5pt, -{Stealth[length=5pt]}] (0,0) -- #2;}
\newcommand\bkvijit[2][0.24]{\makebox[#1\linewidth][l]{#2}}% 章末問題の小問を、幅をそろえて横に並べる（字の間はのびない）
\makeatother

\bkchapter{6}{9}{複素平面：\\かけ算は回転と拡大}{\bklevel{4}}{複素平面：かけ算は回転と拡大}
\bkmichishirube{4-2,4-4,4-7,4-9,4-15,6-8}{6-10,6-13,7-6}

\begin{goal}
\begin{itemize}[label={\color{\bkTT}●}, leftmargin=1.6\zw, itemsep=0pt, topsep=1pt]
\item \bklevel{4}\enspace 複素数を平面の点で表し、絶対値と偏角を求められる。
\item \bklevel{4}\enspace かけ算が「大きさはかけ算、角はたし算」になることを、図と式で説明できる。
\item \bklevel{4}\enspace 複素数のかけ算で、点を回したり、ロボットのうでの先の位置を求めたりできる。
\end{itemize}
\end{goal}

\begin{hatena}
\noindent\begin{minipage}[c]{0.6\linewidth}
\raggedright
6-8で、$j$ をかけると、点が原点のまわりに90°回ることを見ました。では、$1+j$ をかけると、どうなるでしょう。

1に $1+j$ をかけると $1+j$。もう一度かけると $(1+j)^2=2j$。点はだんだん回りながら、原点から遠ざかっていくようです。何回かけると、点はもとの向き（右向き）にもどるでしょうか。
\end{minipage}\hfill
\begin{minipage}[c]{0.37\linewidth}
\centering
\begin{tikzpicture}[x=9mm, y=9mm]
  \draw[bkGray, line width=0.6pt, -{Stealth[length=3.5pt]}] (-2.6,0) -- (1.8,0);
  \draw[bkGray, line width=0.6pt, -{Stealth[length=3.5pt]}] (0,-0.5) -- (0,2.7);
  \draw[bkBlue, line width=1.2pt, -{Stealth[length=4pt]}] (0,0) -- (1,0);
  \draw[bkOrange, line width=1.2pt, -{Stealth[length=4pt]}] (0,0) -- (1,1);
  \draw[bkPurple, line width=1.2pt, -{Stealth[length=4pt]}] (0,0) -- (0,2);
  \draw[bkGray, line width=0.7pt, dash pattern=on 2pt off 1.5pt, -{Stealth[length=3.5pt]}] (1.05,0.15) -- (1.02,0.85);
  \draw[bkGray, line width=0.7pt, dash pattern=on 2pt off 1.5pt, -{Stealth[length=3.5pt]}] (0.85,1.15) -- (0.15,1.92);
  \draw[bkGray, line width=0.7pt, dash pattern=on 2pt off 1.5pt, -{Stealth[length=3.5pt]}] (-0.15,2.0) -- (-1.6,1.6);
  \node[font=\sffamily\footnotesize, anchor=north, inner sep=1pt] at (1,-0.12) {1};
  \node[font=\sffamily\footnotesize, anchor=west, inner sep=1pt] at (1.12,1.05) {$1+j$};
  \node[font=\sffamily\footnotesize, anchor=west, inner sep=1pt] at (0.12,2.25) {$2j$};
  \node[font=\sffamily\bfseries\small, text=bkOrange!80!black] at (-2.0,1.55) {？};
\end{tikzpicture}
\end{minipage}
\end{hatena}

\noindent 6-8では、複素数 $a+bj$ を平面の点 $(a,\ b)$ と考えました。この章では、この平面の上で、複素数のたし算・かけ算・わり算が、それぞれどんな動きになるのかを調べます。かけ算は「回転と拡大」という、図形のとてもきれいな動きになります。この性質のおかげで、電気の交流やロボットの計算が、かけ算1回でできるようになるのです。

はじめての考えが続きますが、1回でわからなくても大丈夫です。表計算で点を動かして、目で確かめながら進みます。% Ver2 No.24 声かけ

%% ---------------- やってみよう・思い出そう ----------------
\begin{yattemiyou}[表計算で、$1+j$ を何回もかけよう]
$a+bj$ に $1+j$ をかけると、6-8のかけ算で
\bkvijline{$(a+bj)(1+j)=a+aj+bj+bj^2=(a-b)+(a+b)j$}
になります。新しい実部は $a-b$、新しい虚部は $a+b$ です。これを表計算で何回もくり返します。
\begin{enumerate}[label={\bkcircled[bkGreen]{\arabic*}}, leftmargin=2.4\zw, itemsep=1pt, topsep=3pt]
\item A列に回数 $n$（0〜8）、B列に実部、C列に虚部を入れる。2行めは1と0（はじめの数1）。
\item B3に \texttt{=B2-C2}、C3に \texttt{=B2+C2} と入れて、下へコピーする。
\item D列に原点からの距離 \texttt{=SQRT(B2\^{}2+C2\^{}2)}、E列に向きの角 \texttt{=DEGREES(ATAN2(B2,C2))} を入れて、下へコピーする。
\item 方眼紙に点を打ち、距離と角が、1回ごとにどう変わるかを調べよう。
\end{enumerate}
\begin{bkcenter}
\bksheet{5}{19}{%
1 & $n$ & 実部 & 虚部 & 距離 & 角（度） \\
2 & 0 & 1 & 0 & （③の式） & （③の式） \\
3 & 1 & \texttt{=B2-C2} & \texttt{=B2+C2} & （コピー） & （コピー） \\}
\end{bkcenter}
\end{yattemiyou}

\begin{bkcenter}
\bksheet{5}{18}{%
1 & $n$ & 実部 & 虚部 & 距離 & 角（度） \\
2 & 0 & 1 & 0 & 1 & 0 \\
3 & 1 & 1 & 1 & 1.414 & 45 \\
4 & 2 & 0 & 2 & 2 & 90 \\
5 & 3 & $-2$ & 2 & 2.828 & 135 \\
6 & 4 & $-4$ & 0 & 4 & 180 \\
7 & 5 & $-4$ & $-4$ & 5.657 & $-135$ \\
8 & 6 & 0 & $-8$ & 8 & $-90$ \\
9 & 7 & 8 & $-8$ & 11.314 & $-45$ \\
10 & 8 & 16 & 0 & 16 & 0 \\}
\bktabcaption{計算の結果（距離は小数第3位まで）}
\end{bkcenter}


\begin{bkfigure}
\begin{tikzpicture}[x=4.6mm, y=4.6mm]
  \draw[bkLine, line width=0.4pt, step=1] (-5.5,-8.6) grid (9.5,2.6);
  \draw[bkGray, line width=0.6pt, -{Stealth[length=4pt]}] (-5.7,0) -- (9.9,0);
  \draw[bkGray, line width=0.6pt, -{Stealth[length=4pt]}] (0,-8.8) -- (0,3);
  \foreach \x/\y in {1/0, 1/1, 0/2, -2/2, -4/0, -4/-4, 0/-8, 8/-8} {\draw[bkBlue!55, line width=0.8pt] (0,0) -- (\x,\y);}
  \draw[bkOrange, line width=1pt, dash pattern=on 3pt off 1.5pt, -{Stealth[length=4pt]}] (1,0) -- (1,1) -- (0,2) -- (-2,2) -- (-4,0) -- (-4,-4) -- (0,-8) -- (8,-8);
  \foreach \x/\y in {1/0, 1/1, 0/2, -2/2, -4/0, -4/-4, 0/-8, 8/-8} {\fill[bkBlue] (\x,\y) circle (0.2);}
  \node[font=\sffamily\footnotesize, anchor=north west, fill=white, inner sep=1pt] at (1.15,-0.15) {1（$n=0$）};
  \node[font=\sffamily\footnotesize, anchor=west, fill=white, inner sep=1pt] at (1.3,1.2) {$1+j$};
  \node[font=\sffamily\footnotesize, anchor=west, fill=white, inner sep=1pt] at (0.35,2.0) {$2j$};
  \node[font=\sffamily\footnotesize, anchor=south, fill=white, inner sep=1pt] at (-2,2.3) {$-2+2j$};
  \node[font=\sffamily\footnotesize, anchor=east, fill=white, inner sep=1pt] at (-4.3,0.6) {$-4$};
  \node[font=\sffamily\footnotesize, anchor=east, fill=white, inner sep=1pt] at (-4.3,-4) {$-4-4j$};
  \node[font=\sffamily\footnotesize, anchor=east, fill=white, inner sep=1pt] at (-0.4,-8) {$-8j$};
  \node[font=\sffamily\footnotesize, anchor=west, fill=white, inner sep=1pt] at (8.35,-8) {$8-8j$};
\end{tikzpicture}
\bkfigcaption{1に $1+j$ を1回ずつかけた点（$n=0$〜7）。点は45°ずつ左回りに回りながら、原点から遠ざかる。$n=8$ の16は、図の右の外。}\label{bk:6-9:fig-uzu}
\end{bkfigure}

\noindent 表のD列を見ると、距離は $1\to1.414\to2\to2.828\to4\to\cdots$ と、1回ごとに約1.414倍になっています。1.414は $\sqrt2$ です（2-8）。E列の角は、1回ごとに45°ずつふえています（$135°$ の次の $180°$ をこえた所は、$-135°$ と表示されます。$-135°$ は、左回りに $225°$ 回ったのと同じ向きです。4-8）。そして、$1+j$ という数そのものの、原点からの距離は $\sqrt{1^2+1^2}=\sqrt2$、向きの角は45°です。$1+j$ をかけるたびに、「$1+j$ の距離の $\sqrt2$ 倍に拡大し、$1+j$ の角の45°だけ回す」ことが起きているようです。この章で、そのわけを調べます。

\begin{omoidasou}
\bkrecall{6-8}{$j^2=-1$。$(a+bj)(c+dj)=(ac-bd)+(ad+bc)j$。$j$ をかけると左回りに90°回る。\\$z=a+bj$ と共役な数 $\bar z=a-bj$。$z\bar z=a^2+b^2$}
\bkrecall{4-4}{原点から点 $(a,\ b)$ までの距離は $\sqrt{a^2+b^2}$（三平方の定理）。}
\bkrecall{4-15}{極座標 $(r,\ \theta)$：$x=r\cos\theta$、$y=r\sin\theta$。$\tan^{-1}$ は右半分の角しか答えないので、点の区画を見て直す。}
\bkrecall{4-9}{加法定理（単位円の上の点への道すじを、まるごと回して作った）：\\$\cos(a+b)=\cos a\cos b-\sin a\sin b$\\$\sin(a+b)=\sin a\cos b+\cos a\sin b$}
\bkrecall{4-11}{ベクトルのたし算・ひき算は、成分ごと。2点の間の矢印は、\mbox{（終わりの点）${}-{}$（始めの点）}。}
\bkrecall{4-7}{2リンクのうでの手の位置：$x=L_1\cos\theta_1+L_2\cos(\theta_1+\theta_2)$、$y=L_1\sin\theta_1+L_2\sin(\theta_1+\theta_2)$}
\end{omoidasou}

%% ---------------- 数学の道具 ----------------
\bksection{doguu}{数学の道具}

\bksubsection{1}{複素数を点で表す：複素平面}

\leavevmode\label{bk:6-9:heimen}6-8では、複素数 $a+bj$ を、平面の点 $(a,\ b)$ に置きました。この平面を\textbf{複素平面}（ふくそへいめん）といいます。\bkidx{ふくそへいめん}{複素平面}横の軸は実部を表すので\textbf{実軸}、たての軸は虚部を表すので\textbf{虚軸}といいます。実軸は、これまでの数直線そのものです。

\begin{bkfigure}
\begin{tikzpicture}[x=10mm, y=10mm]
  \draw[bkLine, line width=0.4pt, step=1] (-2.6,-1.6) grid (4.6,3.6);
  \bkvijaxes{-2.8}{5.2}{-1.8}{4.1}
  \foreach \k in {-2,-1,1,2,3,4} {\node[font=\sffamily\footnotesize, anchor=north, inner sep=1pt, fill=white] at (\k,-0.1) {$\k$};}
  \foreach \k/\l in {1/j,2/2j,3/3j} {\node[font=\sffamily\footnotesize, anchor=east, inner sep=1pt, fill=white] at (-0.1,\k) {$\l$};}
  \node[font=\sffamily\footnotesize, anchor=east, inner sep=1pt, fill=white] at (-0.1,-1) {$-j$};
  \draw[bkBlue, line width=0.7pt, dash pattern=on 2pt off 1.5pt] (3,0) -- (3,2) (0,2) -- (3,2);
  \bkvijarrow{bkBlue}{(3,2)}
  \fill[bkBlue] (3,2) circle (0.09);
  \node[font=\sffamily\footnotesize, anchor=west, fill=white, inner sep=1pt] at (3.15,2.2) {$z=3+2j$};
  \fill[bkOrange] (-2,1) circle (0.09);
  \node[font=\sffamily\footnotesize, anchor=south, fill=white, inner sep=1pt] at (-2,1.15) {$-2+j$};
  \fill[bkPurple] (1,-1) circle (0.09);
  \node[font=\sffamily\footnotesize, anchor=west, fill=white, inner sep=1pt] at (1.15,-1.05) {$1-j$};
  \node[font=\sffamily\footnotesize, text=bkBlue!80!black, anchor=south, fill=white, inner sep=1pt] at (2.1,0.12) {実部3};
  \node[font=\sffamily\footnotesize, text=bkBlue!80!black, anchor=west, fill=white, inner sep=1pt] at (3.1,1.0) {虚部2};
\end{tikzpicture}
\bkfigcaption{複素平面。$z=3+2j$ は、実部3・虚部2の点。原点からの矢印でも表す。}
\end{bkfigure}

複素数は、原点からその点への矢印（ベクトル。4-10）で表すこともできます。$z=3+2j$ は「右へ3、上へ2」の矢印です。点で考えるか、矢印で考えるかは、場面によって使い分けます。

複素数を平面の点で表す考えは、ノルウェー生まれの測量技師ウェッセル（1797年に発表）と、パリのアルガン（1806年）が、別々に思いつきました。ガウスが1831年にこの考えを広めたので、複素平面を「ガウス平面」ともいいます。

\bksubsection{2}{たし算・ひき算は、矢印のたし算・ひき算}

\leavevmode\label{bk:6-9:tasu}複素数のたし算は、実部どうし、虚部どうしをたしました（6-8⑤）。これは、4-11のベクトルの成分ごとのたし算と同じです。だから、複素平面の上では、\textbf{たし算は矢印をつなぐこと}です。

\begin{bkfigure}
\begin{tikzpicture}[x=8mm, y=8mm]
  \begin{scope}
    \draw[bkLine, line width=0.4pt, step=1] (-0.6,-0.6) grid (4.6,3.6);
    \draw[bkGray, line width=0.6pt, -{Stealth[length=4pt]}] (-0.8,0) -- (5,0);
    \draw[bkGray, line width=0.6pt, -{Stealth[length=4pt]}] (0,-0.8) -- (0,4);
    \draw[bkOrange!60, line width=0.9pt, dash pattern=on 2.5pt off 1.5pt, -{Stealth[length=4pt]}] (3,1) -- (4,3);
    \draw[bkBlue!60, line width=0.9pt, dash pattern=on 2.5pt off 1.5pt, -{Stealth[length=4pt]}] (1,2) -- (4,3);
    \bkvijarrow{bkBlue}{(3,1)}
    \bkvijarrow{bkOrange}{(1,2)}
    \bkvijarrow{bkPurple}{(4,3)}
    \node[font=\sffamily\footnotesize, anchor=north west, fill=white, inner sep=1pt] at (3.05,0.95) {$z=3+j$};
    \node[font=\sffamily\footnotesize, anchor=south east, fill=white, inner sep=1pt] at (0.95,2.1) {$w=1+2j$};
    \node[font=\sffamily\footnotesize, anchor=south, fill=white, inner sep=1pt] at (4,3.15) {$z+w=4+3j$};
  \end{scope}
  \begin{scope}[xshift=62mm]
    \draw[bkLine, line width=0.4pt, step=1] (-0.6,-0.6) grid (4.6,5.6);
    \draw[bkGray, line width=0.6pt, -{Stealth[length=4pt]}] (-0.8,0) -- (5,0);
    \draw[bkGray, line width=0.6pt, -{Stealth[length=4pt]}] (0,-0.8) -- (0,6);
    \draw[bkPurple!60, line width=0.9pt, dash pattern=on 2.5pt off 1.5pt, -{Stealth[length=4pt]}] (0,0) -- (3,4);
    \draw[bkPurple, line width=1.5pt, -{Stealth[length=5pt]}] (1,1) -- (4,5);
    \fill[bkBlue] (1,1) circle (0.1);
    \fill[bkOrange] (4,5) circle (0.1);
    \node[font=\sffamily\footnotesize, anchor=north west, fill=white, inner sep=1pt] at (1.1,0.95) {$z_1=1+j$};
    \node[font=\sffamily\footnotesize, anchor=south, fill=white, inner sep=1pt] at (4,5.15) {$z_2=4+5j$};
    \node[font=\sffamily\footnotesize, anchor=east, fill=white, inner sep=1pt, align=right] at (2.3,3.4) {$z_2-z_1=3+4j$\\長さ5};
  \end{scope}
\end{tikzpicture}
\bkfigcaption{左：$z+w$ は、$z$ の矢印の先に $w$ の矢印をつないだ所（平行四辺形の対角線）。右：$z_2-z_1$ は、$z_1$ から $z_2$ への矢印。原点から引きなおすと（点線）$3+4j$。}\label{bk:6-9:fig-tasu}
\end{bkfigure}

\bkvijhead{2つの点の距離}
ひき算 $z_2-z_1$ は、「$z_1$ に何をたすと $z_2$ になるか」なので、$z_1$ から $z_2$ へ向かう矢印です（図\ref{bk:6-9:fig-tasu}の右）。だから、2つの点 $z_1$、$z_2$ の距離は、$z_2-z_1$ の矢印の長さです。$z_1=1+j$、$z_2=4+5j$ なら、$z_2-z_1=3+4j$ で、長さは $\sqrt{3^2+4^2}=\sqrt{25}=5$ です。

\bkvijhead{実数をかけると}
実数 $k$ をかけた $kz$ は、4-10のベクトルの実数倍と同じで、$k>0$ なら同じ向きに $k$ 倍、$k<0$ なら反対向きに $|k|$ 倍の矢印です。$-z$ は、$z$ を原点のまわりに180°回した点です（6-8の $\times(-1)$）。

\bksubsection{3}{原点からの距離：絶対値}

\leavevmode\label{bk:6-9:zettai}複素数 $z=a+bj$ の点と原点との距離を、$z$ の\textbf{絶対値}（ぜったいち）といい、$|z|$ と書きます。\bkidx{ぜったいち}{絶対値（複素数の）}三平方の定理（4-4）で、
\bkvijframe[80mm]{\large $|a+bj|=\sqrt{a^2+b^2}$}
たとえば $|3+4j|=\sqrt{9+16}=5$、$|1+j|=\sqrt2$、$|-2j|=\sqrt{0+4}=2$ です。実数 $a$（虚部0）なら $|a|=\sqrt{a^2}$ で、3-2で学んだ「0からの距離」としての絶対値と同じです。$|-3|=3$。

$\sqrt{\phantom{2}}$ の中に入れるのは、実部と虚部（どちらも実数）の2乗です。$j$ は入れません。また、6-8⑦の $z\bar z=a^2+b^2$ から、
\bkvijline{$|z|^2=z\bar z$}
です。共役な $\bar z$ は、$z$ を実軸で裏返した点なので、原点からの距離は同じで、$|\bar z|=|z|$ です。6-8のもっと知りたいで見たように、複素数には大小がありません。複素数の「大きさ」をくらべるときは、絶対値をくらべます。

\bksubsection{4}{向きの角：偏角と極形式}

\leavevmode\label{bk:6-9:kyoku}原点から $z$ への矢印が、実軸の正の向きから左回りに何度回った向きか、その角を $z$ の\textbf{偏角}（へんかく）といい、$\arg z$（アーギュメント・ゼット）と書きます。\bkidx{へんかく}{偏角}4-15の極座標の $\theta$ と同じものです。

\begin{bkfigure}
\begin{tikzpicture}[x=1mm, y=1mm]
  \begin{scope}
    \draw[bkGray, line width=0.6pt, -{Stealth[length=4pt]}] (-6,0) -- (46,0) node[anchor=north east, font=\sffamily\footnotesize, inner sep=1pt, yshift=-1mm] {実軸};
    \draw[bkGray, line width=0.6pt, -{Stealth[length=4pt]}] (0,-5) -- (0,32) node[anchor=north west, font=\sffamily\footnotesize, inner sep=1pt, xshift=1mm] {虚軸};
    \draw[bkBlue, line width=1.6pt] (0,0) -- ({36*cos(35)},0);
    \draw[bkOrange, line width=1.6pt] ({36*cos(35)},0) -- (35:36);
    \draw[bkInk, line width=1.1pt, -{Stealth[length=5pt]}] (0,0) -- (35:36);
    \fill[bkInk] (35:36) circle (1);
    \draw[bkInk, line width=0.6pt] (9,0) arc (0:35:9);
    \node[font=\sffamily\footnotesize] at (17:13.5) {$\theta$};
    \node[font=\sffamily\footnotesize, anchor=south east, inner sep=1pt] at (35:19) {$r=|z|$};
    \node[font=\sffamily\footnotesize, text=bkBlue!80!black, anchor=north, inner sep=1pt] at (15,-1.2) {$a=r\cos\theta$};
    \node[font=\sffamily\footnotesize, text=bkOrange!80!black, anchor=west, inner sep=1pt] at ({36*cos(35)+1},10) {$b=r\sin\theta$};
    \node[font=\sffamily\footnotesize, anchor=south west, inner sep=1pt] at ($(35:36)+(0.8,0.5)$) {$z=a+bj$};
  \end{scope}
  \begin{scope}[xshift=90mm, yshift=12mm, x=11mm, y=11mm]
    \draw[bkGray, line width=0.6pt, -{Stealth[length=4pt]}] (-2,0) -- (2,0);
    \draw[bkGray, line width=0.6pt, -{Stealth[length=4pt]}] (0,-1.6) -- (0,2);
    \draw[bkBlue, line width=1.3pt, -{Stealth[length=4pt]}] (0,0) -- (1,1);
    \draw[bkOrange, line width=1.3pt, -{Stealth[length=4pt]}] (0,0) -- (-1,1.732);
    \draw[bkPurple, line width=1.3pt, -{Stealth[length=4pt]}] (0,0) -- (-1,-1);
    \draw[bkBlue, line width=0.6pt] (0.42,0) arc (0:45:0.42);
    \draw[bkOrange, line width=0.6pt] (0.62,0) arc (0:120:0.62);
    \draw[bkPurple, line width=0.6pt, -{Stealth[length=3pt]}] (0.82,0) arc (0:-135:0.82);
    \node[font=\sffamily\footnotesize, anchor=west, inner sep=1pt] at (1.05,1.0) {$1+j$：45°};
    \node[font=\sffamily\footnotesize, anchor=south, inner sep=1pt, fill=white] at (-1,1.8) {$-1+\sqrt3j$：120°};
    \node[font=\sffamily\footnotesize, anchor=north, inner sep=1pt, fill=white] at (-1,-1.08) {$-1-j$：$-135°$};
  \end{scope}
\end{tikzpicture}
\bkfigcaption{左：絶対値 $r$ と偏角 $\theta$。右：3つの複素数の偏角。$-1-j$ は、右回りに135°（$-135°$）。}\label{bk:6-9:fig-henkaku}
\end{bkfigure}

\bkvijhead{偏角の範囲}
角は、360°回るともとの向きにもどるので（4-8の一般角）、$-135°$ と $225°$ は同じ向きです。この本では、偏角を $-180°$ より大きく $180°$ 以下の角で答えることにします。Python の \texttt{cmath.phase}、表計算の \texttt{IMARGUMENT} や \texttt{ATAN2} が答える角と同じ範囲です（4-15では $0°$〜$360°$ で答えましたが、どちらも同じ向きを表しています）。実軸の正の向きにある点（正の実数）は $0°$、負の実数は $180°$ です。原点（0）だけは、向きが決まらないので、偏角を考えません。

\bkvijhead{偏角の求め方：区画に注意}
$\tan\theta=\dfrac{b}{a}$（傾き。4-6）なので、電卓の $\tan^{-1}$ が使えそうです。ところが4-15で学んだように、$\tan^{-1}$ は右半分（$-90°$〜$90°$）の角しか答えません。たとえば $-1-j$ では $\tan\theta=\dfrac{-1}{-1}=1$ で、$\tan^{-1}1=45°$ と出ますが、$-1-j$ は左下にあるので、45°は反対向きです。だから、点がどの区画にあるかを見て直します。
\begin{bkcenter}
\begin{bktable}{colspec={Q[l,m,wd=34mm] Q[l,m,wd=76mm]}, rowsep=1.5pt}
点のある所 & 偏角の求め方（$-180°<\theta\leqq180°$） \\
右半分（$a>0$） & $\tan^{-1}\dfrac{b}{a}$ のまま \\
左上（$a<0$、$b\geqq0$） & $\tan^{-1}\dfrac{b}{a}$ に180°をたす \\
左下（$a<0$、$b<0$） & $\tan^{-1}\dfrac{b}{a}$ から180°をひく \\
虚軸の上（$a=0$） & $b>0$ なら90°、$b<0$ なら $-90°$ \\
\end{bktable}
\bktabcaption{いちばん確かなのは、先に点をかいて、だいたいの向きを見ること。}
\end{bkcenter}
$-1+\sqrt3\,j$ なら、$\tan^{-1}\dfrac{\sqrt3}{-1}=\tan^{-1}(-\sqrt3)=-60°$ で、左上なので180°をたして $120°$。$-1-j$ なら、$\tan^{-1}1=45°$ から180°をひいて $-135°$ です（図\ref{bk:6-9:fig-henkaku}の右）。

\bkvijhead{極形式}
絶対値 $r$ と偏角 $\theta$ がわかると、実部と虚部は、4-15の極座標と同じで、$a=r\cos\theta$、$b=r\sin\theta$ です（図\ref{bk:6-9:fig-henkaku}の左）。だから、
\bkvijframe[100mm]{\large $z=a+bj=r(\cos\theta+j\sin\theta)$}
この書き方を、複素数の\textbf{極形式}（きょくけいしき）といいます。\bkidx{きょくけいしき}{極形式}「距離と向き」で書いた複素数です。たとえば $1+j$ は、$r=\sqrt2$、$\theta=45°$ なので $\sqrt2(\cos45°+j\sin45°)$。反対に、$2(\cos120°+j\sin120°)$ は、$\cos120°=-\dfrac12$、$\sin120°=\dfrac{\sqrt3}{2}$ なので、
\bkvijline{$2\Bigl(-\dfrac{1}{2}+\dfrac{\sqrt3}{2}j\Bigr)=-1+\sqrt3\,j$}
です。かっこの中の $\cos\theta+j\sin\theta$ は、絶対値が $\sqrt{\cos^2\theta+\sin^2\theta}=1$（4-8）の、単位円の上の点です。極形式は「単位円の上の、角 $\theta$ の点を、$r$ 倍した点」と読めます。

\bksubsection{5}{かけ算は「大きさはかけ算、角はたし算」}

\leavevmode\label{bk:6-9:kakeru}やってみようでは、$1+j$ をかけるたびに、距離が $\sqrt2$ 倍、角が45°ずつふえました。これが $1+j$ だけでなく、どんな複素数のかけ算でも成り立つことを、3通りで確かめます。

\bkvijhead{見方1：数で確かめる}
$z=1+j$（絶対値 $\sqrt2$、偏角45°）と、$w=1+\sqrt3\,j$（絶対値 $\sqrt{1+3}=2$、偏角60°）をかけます。
\begin{bkhenkei}
zw=(1+j)(1+\sqrt3\,j) & 問題の式 \\
=1+\sqrt3\,j+j+\sqrt3\,j^2 & 展開する（6-8⑥） \\
=(1-\sqrt3)+(1+\sqrt3)j & $j^2=-1$。実部と虚部を集める \\
\approx-0.732+2.732j & $\sqrt3\approx1.732$ \\
\end{bkhenkei}
この答えの絶対値と偏角を求めます。
\begin{itemize}[leftmargin=1.2\zw, itemsep=2pt, topsep=2pt, label={\color{\bkTT}●}]
\item 絶対値：$\sqrt{(-0.732)^2+2.732^2}=\sqrt{0.536+7.464}=\sqrt8\approx2.828$。これは $\sqrt2\times2=2\sqrt2$（$\sqrt8=\sqrt{4\times2}=2\sqrt2$）。
\item 偏角：左上の点なので、$\tan^{-1}\dfrac{2.732}{-0.732}=\tan^{-1}(-3.732)=-75°$ に180°をたして $105°$。これは $45°+60°$。
\end{itemize}
絶対値は $\sqrt2$ と2の\textbf{かけ算}、偏角は45°と60°の\textbf{たし算}になりました。

\bkvijhead{見方2：図で考える（道すじごと回す）}
$w=c+dj$ を $z$ にかけると、分配法則で
\bkvijline{$zw=z(c+dj)=c\times z+d\times(jz)$}
です。$c\times z$ は「$z$ の矢印を $c$ 倍」、$jz$ は「$z$ を90°回した矢印」（6-8）なので、$d\times(jz)$ は「$z$ を90°回して $d$ 倍」です。$z=2+j$、$w=1+j$（$c=1$、$d=1$）で図にします。

\begin{bkfigure}
\begin{tikzpicture}[x=11mm, y=11mm]
  \begin{scope}
    \draw[bkGray, line width=0.6pt, -{Stealth[length=4pt]}] (-0.4,0) -- (1.9,0);
    \draw[bkGray, line width=0.6pt, -{Stealth[length=4pt]}] (0,-0.4) -- (0,1.7);
    \fill[bkBlue!10] (0,0) -- (1,0) -- (1,1) -- cycle;
    \draw[bkBlue, line width=1.5pt, -{Stealth[length=5pt]}] (0,0) -- (1,0);
    \draw[bkOrange, line width=1.5pt, -{Stealth[length=5pt]}] (1,0) -- (1,1);
    \draw[bkPurple, line width=1.5pt, -{Stealth[length=5pt]}] (0,0) -- (1,1);
    \draw[bkInk, line width=0.5pt] (0.88,0) -- (0.88,0.12) -- (1,0.12);
    \node[font=\sffamily\footnotesize, anchor=north, inner sep=1pt] at (0.5,-0.08) {1};
    \node[font=\sffamily\footnotesize, anchor=west, inner sep=1pt] at (1.06,0.5) {$j$};
    \node[font=\sffamily\footnotesize, anchor=south east, inner sep=1pt] at (0.42,0.78) {$w=1+j$};
    \node[font=\sffamily\small, anchor=north] at (0.75,-0.5) {$w$ への道すじ};
  \end{scope}
  \begin{scope}[xshift=40mm]
    \draw[bkLine, line width=0.4pt, step=1] (-0.5,-0.5) grid (2.5,3.5);
    \draw[bkGray, line width=0.6pt, -{Stealth[length=4pt]}] (-0.6,0) -- (2.9,0);
    \draw[bkGray, line width=0.6pt, -{Stealth[length=4pt]}] (0,-0.6) -- (0,3.8);
    \fill[bkBlue!10] (0,0) -- (2,1) -- (1,3) -- cycle;
    \draw[bkBlue, line width=1.5pt, -{Stealth[length=5pt]}] (0,0) -- (2,1);
    \draw[bkOrange, line width=1.5pt, -{Stealth[length=5pt]}] (2,1) -- (1,3);
    \draw[bkPurple, line width=1.5pt, -{Stealth[length=5pt]}] (0,0) -- (1,3);
    \draw[bkInk, line width=0.5pt] ($(2,1)+(-0.107,-0.054)$) -- ++(-0.054,0.107) -- ++(0.107,0.054);
    \draw[bkInk, line width=0.6pt] (0.7,0) arc (0:26.57:0.7);
    \node[font=\sffamily\footnotesize] at (13:1.25) {26.6°};
    \node[font=\sffamily\footnotesize, anchor=north west, inner sep=1pt, fill=white] at (2.05,0.95) {$z=2+j$};
    \node[font=\sffamily\footnotesize, anchor=west, inner sep=1pt, fill=white] at (1.6,2.15) {$jz=-1+2j$};
    \node[font=\sffamily\footnotesize, anchor=south, inner sep=1pt, fill=white] at (1,3.08) {$zw=1+3j$};
  \end{scope}
  \node[font=\sffamily\small, anchor=north west, align=left, text width=44mm] at (82mm,36mm) {%
  左の三角形（1と $j$ の道すじ）を、\\
  \textcolor{bkBlue!80!black}{・$|z|=\sqrt5$ 倍に拡大して}\\
  \textcolor{bkBlue!80!black}{・$z$ の偏角26.6°だけ回す}\\
  と、右の三角形になる。\\[4pt]
  だから $zw$ の矢印は、\\
  $w$ の矢印を $\sqrt5$ 倍にして、\\26.6°回したもの。};
\end{tikzpicture}
\bkfigcaption{$w=1\times1+1\times j$ への道すじ（左）と、$zw=1\times z+1\times(jz)$ への道すじ（右）。直角三角形の形は同じで、大きさと向きだけがちがう。}\label{bk:6-9:fig-michi}
\end{bkfigure}

右の三角形の2つの辺は、$z$ の矢印（長さ $|z|$）と、それを90°回した $jz$ の矢印（長さ $|z|$）です。左の三角形の辺1と $j$ の両方を、それぞれ $|z|$ 倍にして $z$ の向きに回したものなので、三角形の形はそのままで、$|z|$ 倍に拡大されて、$z$ の偏角だけ回っています（4-3の拡大）。だから、斜辺の $zw$ は、$w$ を $|z|$ 倍に拡大して、$\arg z$ だけ回した矢印です。

これは、4-9で加法定理を作ったときの「道すじを、まるごと回す」図と同じ考えです。$c$、$d$ がほかの数でも、三角形の横とたてが $c$ と $d$ になるだけで、同じことが言えます。

\bkvijhead{見方3：式で確かめる（加法定理）}
$z=r_1(\cos\alpha+j\sin\alpha)$、$w=r_2(\cos\beta+j\sin\beta)$ とします（$\alpha$ はアルファ、$\beta$ はベータ。3-10）。
実数の $r_1$、$r_2$ を前に集めると、
\bkvijline{$zw=r_1r_2(\cos\alpha+j\sin\alpha)(\cos\beta+j\sin\beta)$}
です。ここから1行ずつ進めます。
\begin{bkhenkei}
=r_1r_2\{\cos\alpha\cos\beta+j\cos\alpha\sin\beta & 展開する（4つの組み合わせ） \\
\quad+j\sin\alpha\cos\beta+j^2\sin\alpha\sin\beta\} & \\
=r_1r_2\{(\cos\alpha\cos\beta-\sin\alpha\sin\beta) & $j^2=-1$。$j$ のない項と \\
\quad+j(\sin\alpha\cos\beta+\cos\alpha\sin\beta)\} & $j$ のつく項に分ける \\
=r_1r_2\{\cos(\alpha+\beta)+j\sin(\alpha+\beta)\} & 加法定理（4-9） \\
\end{bkhenkei}
最後の式は、絶対値が $r_1r_2$、偏角が $\alpha+\beta$ の複素数の極形式です。3通りとも、同じ結論になりました。
\bkvijframe[110mm]{\large $|zw|=|z|\,|w|$\qquad $\arg(zw)=\arg z+\arg w$\\[3pt] かけ算は、大きさはかけ算、角はたし算}
4-9では、加法定理を「道すじを回す」図から作りました。複素数で考えると、加法定理は「かけ算すると角がたされる」ことを、実部と虚部に分けて書いたものだったのです。

\bkvijhead{\bkhatenaref の答え}
$1+j$ は、絶対値 $\sqrt2$、偏角45°です。だから $1+j$ をかけるたびに、点は $\sqrt2$ 倍に拡大されて、45°ずつ左回りに回ります。もとの向き（偏角0°）にもどるのは、$45°\times8=360°$ で、\textbf{8回かけたとき}です。そのとき、距離は $(\sqrt2)^8=2^4=16$ 倍になっているので、$(1+j)^8=16$。やってみようの表の $n=8$ の行と、ぴったり合います。

\bksubsection{6}{特別な数をかける}

\leavevmode\label{bk:6-9:tokubetsu}⑤のきまりで、これまでのかけ算を見なおします。
\begin{bkcenter}
\begin{bktable}{colspec={Q[c,m,wd=28mm] Q[c,m,wd=18mm] Q[c,m,wd=16mm] Q[l,m,wd=48mm]}, rowsep=1.5pt}
かける数 & 絶対値 & 偏角 & 点の動き \\
$j$ & 1 & 90° & 左回りに90°回す（6-8） \\
$-1$ & 1 & 180° & 180°回す（2-2） \\
$-j$ & 1 & $-90°$ & 右回りに90°回す \\
正の実数 $k$ & $k$ & 0° & 回さずに $k$ 倍 \\
$\cos\theta+j\sin\theta$ & 1 & $\theta$ & 大きさを変えずに $\theta$ 回す \\
\end{bktable}
\bktabcaption{絶対値が1の数をかけると、回すだけ。偏角が0°の数をかけると、拡大・縮小だけ。}
\end{bkcenter}
$j\times j$ は、偏角が $90°+90°=180°$、絶対値が $1\times1=1$ なので $-1$。6-8で図から考えた $j^2=-1$ が、⑤のきまりからも出てきます。

\bkvijhead{2回かけると $j$ になる数}
⑤のきまりを使うと、新しい数も見つけられます。2回かけて $j$（絶対値1、偏角90°）になる数は、絶対値が1で、偏角が90°の半分の45°の数です。
\bkvijline{$\cos45°+j\sin45°=\dfrac{\sqrt2}{2}+\dfrac{\sqrt2}{2}j=\dfrac{\sqrt2}{2}(1+j)$}
確かめると、$\Bigl(\dfrac{\sqrt2}{2}\Bigr)^2(1+j)^2=\dfrac{2}{4}\times2j=j$。たしかに $j$ になります。

\bkvijhead{共役な数とのかけ算}
$\bar z$ は、$z$ を実軸で裏返した点なので、絶対値は同じで、偏角は $-\theta$ です。だから $z\bar z$ は、偏角が $\theta+(-\theta)=0°$、絶対値が $|z|^2$ の数、つまり正の実数 $|z|^2$ です。6-8⑦の $z\bar z=a^2+b^2$ が、図の上でも説明できました。

\bksubsection{7}{わり算は「大きさはわり算、角はひき算」}

\leavevmode\label{bk:6-9:waru}わり算はかけ算の逆です（1-14）。$\dfrac{z}{w}=q$ とすると、$z=qw$ です。⑤から $|z|=|q|\,|w|$、$\arg z=\arg q+\arg w$ なので、
\bkvijframe[110mm]{\large $\left|\dfrac{z}{w}\right|=\dfrac{|z|}{|w|}$\qquad $\arg\dfrac{z}{w}=\arg z-\arg w$}
たとえば $\dfrac{1+\sqrt3\,j}{1+j}$ は、絶対値が $\dfrac{2}{\sqrt2}=\sqrt2$、偏角が $60°-45°=15°$ です。6-8⑧の方法で確かめます。
\begin{bkhenkei}
\frac{1+\sqrt3\,j}{1+j}=\frac{(1+\sqrt3\,j)(1-j)}{(1+j)(1-j)} & 分母と共役な $1-j$ をかける \\
=\frac{1-j+\sqrt3\,j-\sqrt3\,j^2}{2} & 展開する。分母は $1^2+1^2=2$ \\
=\frac{(1+\sqrt3)+(\sqrt3-1)j}{2} & $j^2=-1$ \\
\approx1.366+0.366j & $\sqrt3\approx1.732$ \\
\end{bkhenkei}
絶対値は $\sqrt{1.366^2+0.366^2}=\sqrt{1.866+0.134}=\sqrt2$、偏角は右上の点なので $\tan^{-1}\dfrac{0.366}{1.366}=15°$。合っています。

\bksubsection{8}{何回もかける：ド・モアブルの定理}

\leavevmode\label{bk:6-9:demoivre}絶対値 $r$、偏角 $\theta$ の数を $n$ 回かけると、⑤のきまりで、絶対値は $r$ を $n$ 回かけた $r^n$、偏角は $\theta$ を $n$ 回たした $n\theta$ になります。
\bkvijframe[120mm]{\large $\{r(\cos\theta+j\sin\theta)\}^n=r^n(\cos n\theta+j\sin n\theta)$}
$r=1$ のときの式 $(\cos\theta+j\sin\theta)^n=\cos n\theta+j\sin n\theta$ を、\textbf{ド・モアブルの定理}といいます。\bkidx{どもあぶるのていり}{ド・モアブルの定理}フランスに生まれ、イギリスで活やくしたド・モアブル（1667〜1754年）の名前からきています。

\bkvijhead{何回も展開しなくてよい}
$(1+j)^8$ を展開で計算すると、7回もかけ算をしなければなりません。極形式なら、$1+j=\sqrt2(\cos45°+j\sin45°)$ なので、
\begin{bkhenkei}
(1+j)^8 & $1+j$ は絶対値 $\sqrt2$、偏角45° \\
=(\sqrt2)^8(\cos360°+j\sin360°) & ド・モアブル：$45°\times8=360°$ \\
=16\times(1+0j)=16 & $(\sqrt2)^8=(\sqrt2^{\,2})^4=2^4=16$ \\
\end{bkhenkei}
1行で終わりました。

\bkvijhead{$n=2$ のとき：2倍角の式がもう一度出てくる}
$(\cos\theta+j\sin\theta)^2$ を、展開とド・モアブルの2通りで計算します。
\begin{itemize}[leftmargin=1.2\zw, itemsep=2pt, topsep=2pt, label={\color{\bkTT}●}]
\item 展開：$\cos^2\theta+2j\sin\theta\cos\theta+j^2\sin^2\theta=(\cos^2\theta-\sin^2\theta)+j\,(2\sin\theta\cos\theta)$
\item ド・モアブル：$\cos2\theta+j\sin2\theta$
\end{itemize}
2つは同じ数なので、実部どうし、虚部どうしが等しくなり（6-8③）、
\bkvijline{$\cos2\theta=\cos^2\theta-\sin^2\theta$\qquad $\sin2\theta=2\sin\theta\cos\theta$}
4-9の2倍角の式が、複素数のかけ算から出てきました。

\bksubsection{9}{回転の道具：点を回す・うでを動かす}

\leavevmode\label{bk:6-9:kaiten}⑥の表から、\textbf{点を原点のまわりに $\theta$ だけ回すには、$\cos\theta+j\sin\theta$ をかければよい}ことがわかります。絶対値が1なので、大きさは変わりません。

\bkvijhead{点 $(3,\ 1)$ を60°回す}
$\cos60°+j\sin60°=\dfrac12+\dfrac{\sqrt3}{2}j$ をかけます。
\begin{bkhenkei}
(3+j)\Bigl(\frac12+\frac{\sqrt3}{2}j\Bigr) & 問題の式 \\
=\frac32+\frac{3\sqrt3}{2}j+\frac12j+\frac{\sqrt3}{2}j^2 & 展開する \\
=\Bigl(\frac32-\frac{\sqrt3}{2}\Bigr)+\Bigl(\frac{3\sqrt3}{2}+\frac12\Bigr)j & $j^2=-1$。実部と虚部を集める \\
\approx0.634+3.098j & $\sqrt3\approx1.732$ \\
\end{bkhenkei}
回した点は約 $(0.63,\ 3.10)$ です（図\ref{bk:6-9:fig-kaiten}）。

\begin{bkfigure}
\begin{tikzpicture}[x=9mm, y=9mm]
  \draw[bkLine, line width=0.4pt, step=1] (-1.6,-0.6) grid (3.6,3.6);
  \draw[bkGray, line width=0.6pt, -{Stealth[length=4pt]}] (-1.8,0) -- (4,0);
  \draw[bkGray, line width=0.6pt, -{Stealth[length=4pt]}] (0,-0.8) -- (0,4);
  \draw[bkInk!35, line width=0.7pt, dash pattern=on 2pt off 1.5pt] (-15:3.162) arc (-15:120:3.162);
  \bkvijarrow{bkBlue}{(3,1)}
  \bkvijarrow{bkOrange}{(0.634,3.098)}
  \draw[bkInk, line width=0.6pt, -{Stealth[length=3.5pt]}] (18.43:1.4) arc (18.43:78.43:1.4);
  \node[font=\sffamily\footnotesize, fill=white, inner sep=1pt] at (48:1.85) {60°};
  \node[font=\sffamily\footnotesize, anchor=west, fill=white, inner sep=1pt] at (3.1,1.0) {$3+j$};
  \node[font=\sffamily\footnotesize, anchor=south west, fill=white, inner sep=1pt] at (0.72,3.12) {$0.634+3.098j$};
  \node[font=\sffamily\footnotesize, text=bkGray, anchor=west, fill=white, inner sep=1pt] at (-15:3.25) {半径 $\sqrt{10}$ の円};
\end{tikzpicture}
\bkfigcaption{$3+j$ に $\cos60°+j\sin60°$ をかけると、原点のまわりに60°回った点に移る。原点からの距離は変わらない。}\label{bk:6-9:fig-kaiten}
\end{bkfigure}
原点からの距離を確かめると、$\sqrt{0.634^2+3.098^2}\approx3.16$ で、もとの $\sqrt{3^2+1^2}=\sqrt{10}\approx3.16$ と同じです。

文字で書くと、点 $(x,\ y)$ を $\theta$ 回すと、
\bkvijline{$(x+yj)(\cos\theta+j\sin\theta)=(x\cos\theta-y\sin\theta)+(x\sin\theta+y\cos\theta)j$}
で、回した点は $(x\cos\theta-y\sin\theta,\ \ x\sin\theta+y\cos\theta)$ です。$\theta=90°$ なら $(-y,\ x)$ で、4-8の「90°回す」になります。この式は、第7巻の「回転行列」（7-6）でもう一度出てきます。

\bkvijhead{Python で確かめる}
6-8⑩の Python でも、絶対値・偏角・回転が計算できます。\texttt{abs(z)} は絶対値、\texttt{cmath.phase(z)} は偏角（ラジアン。$-\pi$ より大きく $\pi$ 以下）、\texttt{math.degrees} はラジアンを度に直す関数、\texttt{cmath.rect(r, θ)} は絶対値 $r$、偏角 $\theta$（ラジアン）の複素数を作る関数です。

\noindent\begin{minipage}{\linewidth}
\begin{bkcode}[Python]
import cmath, math
z = 3+1j
print(abs(z))
print(math.degrees(cmath.phase(z)))
r = cmath.rect(1, math.radians(60))
print(z * r)
\end{bkcode}
\begin{bkcodeout}
3.1622776601683795
18.43494882292201
(0.6339745962155618+3.098076211353316j)
\end{bkcodeout}
\end{minipage}

\par\vspace{4pt}
1行めは $\sqrt{10}\approx3.162$、2行めは $3+j$ の偏角（$\tan^{-1}\dfrac13\approx18.4°$）、3行めは、さっき手で計算した60°回した点です。

\bkvijhead{ロボットのうでの先の位置}
4-7の2リンクのうでを、複素数で計算します。肩を原点にして、上のうでの長さを $L_1$、下のうでの長さを $L_2$、肩の角を $\theta_1$、ひじの角を $\theta_2$ とします（4-7と同じ決め方）。回転をかけ算で書くために、
\bkvijline{$u_1=\cos\theta_1+j\sin\theta_1$\qquad $u_2=\cos\theta_2+j\sin\theta_2$}
とおきます。$u_1$ をかけると $\theta_1$ 回り、$u_2$ をかけると $\theta_2$ 回ります。

\noindent\begin{minipage}[c]{0.4\linewidth}
\centering
\begin{tikzpicture}[x=0.85mm, y=0.85mm]
  \draw[bkGray, line width=0.6pt, -{Stealth[length=4pt]}] (-6,0) -- (52,0);
  \draw[bkGray, line width=0.6pt, -{Stealth[length=4pt]}] (0,-5) -- (0,58);
  \coordinate (E) at (30:40);
  \coordinate (P) at ($(E)+(90:30)$);
  \draw[bkGray, line width=0.6pt, dash pattern=on 2pt off 1.5pt] (E) -- ($(E)+(30:14)$);
  \draw[bkBlue, line width=3pt, line cap=round] (0,0) -- (E);
  \draw[bkOrange, line width=3pt, line cap=round] (E) -- (P);
  \filldraw[fill=white, draw=bkInk, line width=0.7pt] (0,0) circle (1.6) (E) circle (1.6);
  \fill[bkInk] (P) circle (1.4);
  \draw[bkInk, line width=0.5pt] (11,0) arc (0:30:11);
  \node[font=\sffamily\footnotesize] at (15:15) {$\theta_1$};
  \draw[bkInk, line width=0.5pt] ($(E)+(30:9)$) arc (30:90:9);
  \node[font=\sffamily\footnotesize] at ($(E)+(60:13)$) {$\theta_2$};
  \node[font=\sffamily\footnotesize, text=bkBlue!80!black, anchor=north west, inner sep=1pt] at ($(0,0)!0.55!(E)+(1,-1)$) {$L_1$};
  \node[font=\sffamily\footnotesize, text=bkOrange!80!black, anchor=east, inner sep=1pt] at ($(E)!0.5!(P)+(-2,0)$) {$L_2$};
  \node[font=\sffamily\footnotesize, anchor=south, inner sep=1pt] at ($(P)+(0,2)$) {手};
  \node[font=\sffamily\footnotesize, anchor=north east, inner sep=1pt] at (-1,-1) {肩};
  \node[font=\sffamily\footnotesize, anchor=north west, inner sep=1pt] at ($(E)+(1.5,-1.5)$) {ひじ};
\end{tikzpicture}
\end{minipage}\hfill
\begin{minipage}[c]{0.56\linewidth}
\raggedright
手の位置は、次の2つの動きを順に考えると求められます。
\begin{enumerate}[label={\bkcircled[\bkTF]{\arabic*}}, leftmargin=2.2\zw, itemsep=2pt, topsep=2pt]
\item 上のうでを横にのばしたまま、ひじから先だけを $\theta_2$ 回す。手は $L_1+L_2u_2$ の所。
\item うで全体を、肩のまわりに $\theta_1$ 回す。全体に $u_1$ をかける。
\end{enumerate}
\bkvijline{手$=u_1(L_1+L_2u_2)$}
\end{minipage}

\par\vspace{4pt}
展開すると $L_1u_1+L_2u_1u_2$ で、$u_1u_2$ は偏角 $\theta_1+\theta_2$ の数（⑤）なので、実部と虚部は4-7の式 $x=L_1\cos\theta_1+L_2\cos(\theta_1+\theta_2)$、$y=L_1\sin\theta_1+L_2\sin(\theta_1+\theta_2)$ と同じです。複素数で書くと、「回す」がかけ算1回になるので、関節がいくつあっても、同じ形で順にかけていけば計算できます（\bkreidairef{8}）。

\begin{motto}[4]
\textbf{1の $n$ 乗根：正多角形の頂点}\par
$z^3=1$ になる数 $z$ を考えます。$|z|^3=1$ なので $|z|=1$、偏角は3倍すると360°の整数倍（4-8の一般角）になる角で、$0°$、$120°$、$240°$（$=-120°$）の3つです。
\bkvijline{$z=1$、\quad $\cos120°+j\sin120°=-\dfrac12+\dfrac{\sqrt3}{2}j$、\quad $\cos120°-j\sin120°=-\dfrac12-\dfrac{\sqrt3}{2}j$}
3つの点は、単位円の上の正三角形の頂点です。あとの2つは、6-8の例題8（3）の $s^2+s+1=0$ の解と同じです。$z^3-1=(z-1)(z^2+z+1)$（展開して確かめられる）だからです。同じように、$z^4=1$ の解は $1$、$j$、$-1$、$-j$（正方形）、$z^6=1$ の解は正六角形の頂点です。これらを\textbf{1の $n$ 乗根}といいます。
\par\vspace{2pt}
\begin{center}
\begin{tikzpicture}[x=10mm, y=10mm]
  \foreach \n/\xs in {3/0, 4/36, 6/72} {
    \begin{scope}[xshift=\xs mm]
      \draw[bkGray, line width=0.5pt] (-1.35,0) -- (1.35,0) (0,-1.3) -- (0,1.3);
      \draw[bkInk!35, line width=0.6pt] (0,0) circle (1);
      \draw[bkBlue, line width=1pt] (0:1) \foreach \k in {1,...,\n} {-- ({360*\k/\n}:1)};
      \foreach \k in {1,...,\n} {\fill[bkBlue] ({360*\k/\n}:1) circle (0.07);}
      \node[font=\sffamily\footnotesize, anchor=north] at (0,-1.35) {$z^{\n}=1$};
    \end{scope}}
\end{tikzpicture}
\end{center}
3つの解をたすと、$1+\Bigl(-\dfrac12+\dfrac{\sqrt3}{2}j\Bigr)+\Bigl(-\dfrac12-\dfrac{\sqrt3}{2}j\Bigr)=0$ です。120°ずつずれた3つの矢印をつなぐと、もとにもどるからです。これは、発電所から家庭や工場に送る「三相交流」の3本の電線の電圧の和が0になるしくみと同じです（物理・工学ではこう使う）。
\end{motto}

\begin{senmonsho}[大学・電気の教科書ふうの書き方]
\begin{bkhonbun}
複素数 $z$ は極形式 $z=r(\cos\theta+i\sin\theta)$ で表される。ここで $r=|z|$、$\theta=\arg z$ であり、偏角は $2\pi$ の整数倍を除いて一意に定まる。$|z_1z_2|=|z_1||z_2|$、$\arg(z_1z_2)=\arg z_1+\arg z_2$ が成り立つ。電気工学では、極形式を $\dot V=100\angle30^\circ$ のようにも書く。
\end{bkhonbun}
\bkyaku{数学の本では、角をラジアンで書き、虚数単位を $i$ と書く。「$2\pi$ の整数倍を除いて一意に定まる」は「360°ずつのちがいをのぞけば、ただ1つに決まる」（4-8の一般角）。$100\angle30^\circ$ は「絶対値100、偏角30°」の複素数（読み方は「100、角度30度」など）。電気の本では、複素数で表した電圧や電流に、点をつけて $\dot V$ と書くことがある。}
\end{senmonsho}

\begin{taisetsu}
\def\tsl#1{\makebox[6.2\zw][l]{\sffamily\bfseries\small\color{\bkTT}#1}}
\tsl{複素平面}$a+bj$ は点 $(a,\ b)$。たし算は矢印をつなぐ\\[3pt]
\tsl{絶対値・偏角}$|a+bj|=\sqrt{a^2+b^2}$。偏角は区画を見て決める\\[3pt]
\tsl{極形式}$z=r(\cos\theta+j\sin\theta)$\\[3pt]
\tsl{かけ算}$|zw|=|z|\,|w|$、$\arg(zw)=\arg z+\arg w$（わり算は、わり算とひき算）\\[3pt]
\tsl{ド・モアブル}$(\cos\theta+j\sin\theta)^n=\cos n\theta+j\sin n\theta$\\[3pt]
\tsl{回転}$\theta$ 回すには $\cos\theta+j\sin\theta$ をかける
\end{taisetsu}

\begin{kotoba}
\small
\bkword{複素平面}{複素数 $a+bj$ を点 $(a,\ b)$ で表す平面。横が実軸、たてが虚軸。ガウス平面ともいう。}
\bkword{絶対値}{複素数の点と原点との距離 $|z|$。}
\bkword{偏角}{実軸の正の向きから左回りにはかった、$z$ の向きの角 $\arg z$。}
\bkword{極形式}{複素数を絶対値 $r$ と偏角 $\theta$ で $r(\cos\theta+j\sin\theta)$ と書く書き方。}
\bkword{ド・モアブルの定理}{$(\cos\theta+j\sin\theta)^n=\cos n\theta+j\sin n\theta$。}
\end{kotoba}

\bkkigocards{\bkkigocard{$|z|$}{ぜったいち・ゼット}{複素数 $z$ の点と原点との距離。$|a+bj|=\sqrt{a^2+b^2}$}{absolute value}{6-9}}
{\bkkigocard{$\arg z$}{アーギュメント・ゼット}{偏角。$z$ の向きの角。この本では $-180°<\theta\leqq180°$}{argument}{6-9}}


%% ---------------- 例題 ----------------
\bksection{reidai}{例題でたしかめよう}

\begin{reidai}[絶対値と偏角]
次の複素数の絶対値と偏角を求めましょう（偏角は $-180°<\theta\leqq180°$）。\par
\noindent \bkvijit{（1）$3+4j$}\bkvijit{（2）$-1+\sqrt3\,j$}\bkvijit{（3）$-1-j$}\bkvijit{（4）$-2j$}\par
\kangaekata 絶対値は $\sqrt{a^2+b^2}$（③）。偏角は、まず点の区画を見てから、$\tan^{-1}\dfrac{b}{a}$ を直す（④の表）。\par
\shiki\hangindent=4.3\zw\hangafter=1 （1）$\sqrt{9+16}=5$。右上。$\tan^{-1}\dfrac{4}{3}\approx53.1°$ のまま\par
\bkvijq （2）$\sqrt{1+3}=2$。左上。$\tan^{-1}(-\sqrt3)=-60°$、$-60°+180°=120°$\par
\bkvijq （3）$\sqrt{1+1}=\sqrt2$。左下。$\tan^{-1}1=45°$、$45°-180°=-135°$\par
\bkvijq （4）$\sqrt{0+4}=2$。虚軸の下（$a=0$、$b<0$）で $-90°$\par
\kotae\hangindent=4.3\zw\hangafter=1 （1）5、約53.1°\quad（2）2、120°\quad（3）$\sqrt2$、$-135°$\quad（4）2、$-90°$
\end{reidai}

\begin{reidai}[和・差と、2点の距離]
$z_1=1+j$、$z_2=4+5j$ のとき、$z_1+z_2$、$z_2-z_1$ と、2つの点の距離を求めましょう。\par
\kangaekata たし算・ひき算は成分ごと（②）。2点の距離は、$z_2-z_1$ の絶対値。\par
\shiki\hangindent=2.6\zw\hangafter=1 $z_1+z_2=(1+4)+(1+5)j=5+6j$\par
\bkvija $z_2-z_1=(4-1)+(5-1)j=3+4j$\par
\bkvija 距離は $|3+4j|=\sqrt{9+16}=\sqrt{25}=5$\par
\kotae\hangindent=2.6\zw\hangafter=1 $z_1+z_2=5+6j$、$z_2-z_1=3+4j$、距離5
\end{reidai}

\noindent\begin{minipage}[t]{0.485\linewidth}\vspace{0pt}
\begin{jibun}
次の複素数の絶対値と偏角を求めましょう。\par
（1）$5-12j$\newline（2）$-\sqrt3+j$\par
（3）$-2-2j$\newline（4）$3j$
\bkjans{（1）\nobreak 13、右下で $\tan^{-1}\Bigl(-\dfrac{12}{5}\Bigr)\approx-67.4°$。（2）\nobreak 2、左上で $-30°+180°=150°$。（3）\nobreak $2\sqrt2$、左下で $45°-180°=-135°$。（4）\nobreak 3、90°。}
\end{jibun}
\end{minipage}\hfill
\begin{minipage}[t]{0.485\linewidth}\vspace{0pt}
\begin{jibun}
$z_1=2+j$、$z_2=7+13j$ のとき、$z_1+z_2$、$z_2-z_1$ と、2つの点の距離を求めましょう。
\bkjans{$z_1+z_2=9+14j$、$z_2-z_1=5+12j$、距離 $\sqrt{25+144}=13$。}
\end{jibun}
\end{minipage}

\begin{reidai}[極形式と $a+bj$ の書きかえ]
（1）（2）を $a+bj$ の形に、（3）（4）を極形式に直しましょう。\par
（1）$2(\cos150°+j\sin150°)$\newline（2）$4\{\cos(-60°)+j\sin(-60°)\}$\newline（3）$1+j$\newline（4）$-3$\par
\kangaekata $a=r\cos\theta$、$b=r\sin\theta$（④）。反対は、絶対値と偏角を求めて $r(\cos\theta+j\sin\theta)$ に入れる。\par
\shiki\hangindent=4.3\zw\hangafter=1 （1）$\cos150°=-\dfrac{\sqrt3}{2}$、$\sin150°=\dfrac12$。$2\times\Bigl(-\dfrac{\sqrt3}{2}\Bigr)+2\times\dfrac12j$\par
\bkvijq （2）$\cos(-60°)=\dfrac12$、$\sin(-60°)=-\dfrac{\sqrt3}{2}$。$4\times\dfrac12+4\times\Bigl(-\dfrac{\sqrt3}{2}\Bigr)j$\par
\bkvijq （3）$r=\sqrt2$、右上で $\tan^{-1}1=45°$\par
\bkvijq （4）$r=3$、負の実数なので偏角は180°\par
\kotae\hangindent=4.3\zw\hangafter=1 （1）$-\sqrt3+j$\quad（2）$2-2\sqrt3\,j$\quad（3）$\sqrt2(\cos45°+j\sin45°)$\par
\bkvijq （4）$3(\cos180°+j\sin180°)$
\end{reidai}

\begin{bkfigure}
\begin{tikzpicture}[x=7mm, y=7mm]
  \draw[bkLine, line width=0.4pt, step=1] (-3.6,-3.6) grid (2.6,1.6);
  \draw[bkGray, line width=0.6pt, -{Stealth[length=4pt]}] (-3.8,0) -- (2.9,0);
  \draw[bkGray, line width=0.6pt, -{Stealth[length=4pt]}] (0,-3.8) -- (0,1.9);
  \bkvijarrow{bkBlue}{(-1.732,1)}
  \bkvijarrow{bkOrange}{(2,-3.464)}
  \bkvijarrow{bkPurple}{(1,1)}
  \bkvijarrow{bkGreen!70!black}{(-3,0)}
  \node[font=\sffamily\footnotesize, anchor=south east, fill=white, inner sep=1pt] at (-1.8,1.05) {（1）$-\sqrt3+j$：150°};
  \node[font=\sffamily\footnotesize, anchor=west, fill=white, inner sep=1pt] at (2.1,-3.4) {（2）$2-2\sqrt3\,j$：$-60°$};
  \node[font=\sffamily\footnotesize, anchor=west, fill=white, inner sep=1pt] at (1.1,1.15) {（3）$1+j$：45°};
  \node[font=\sffamily\footnotesize, anchor=north, fill=white, inner sep=1pt] at (-2.6,-0.25) {（4）$-3$：180°};
\end{tikzpicture}
\bkfigcaption{\bkreidairef{3}の4つの数。矢印の長さが絶対値、向きが偏角。}
\end{bkfigure}

\begin{reidai}[極形式のかけ算・わり算]
$z=2(\cos30°+j\sin30°)$、$w=3(\cos60°+j\sin60°)$ のとき、$zw$ と $\dfrac{w}{z}$ を、極形式と $a+bj$ の形の両方で求めましょう。\par
\kangaekata かけ算は、絶対値をかけて偏角をたす（⑤）。わり算は、絶対値をわって偏角をひく（⑦）。\par
\shiki\hangindent=2.6\zw\hangafter=1 $zw$：絶対値 $2\times3=6$、偏角 $30°+60°=90°$\par
\bkvija $6(\cos90°+j\sin90°)=6(0+1j)=6j$\par
\bkvija $\dfrac{w}{z}$：絶対値 $\dfrac32$、偏角 $60°-30°=30°$\par
\bkvija $\dfrac32(\cos30°+j\sin30°)=\dfrac32\Bigl(\dfrac{\sqrt3}{2}+\dfrac12j\Bigr)=\dfrac{3\sqrt3}{4}+\dfrac34j$\par
\kotae\hangindent=2.6\zw\hangafter=1 $zw=6(\cos90°+j\sin90°)=6j$\par
\bkvija $\dfrac{w}{z}=\dfrac32(\cos30°+j\sin30°)=\dfrac{3\sqrt3}{4}+\dfrac34j$
\end{reidai}

\noindent\begin{minipage}[t]{0.485\linewidth}\vspace{0pt}
\begin{jibun}
（1）（2）を $a+bj$ に、（3）（4）を極形式に。\par
（1）$6(\cos30°+j\sin30°)$\par
（2）$2\{\cos(-135°)+j\sin(-135°)\}$\par
（3）$-1+\sqrt3\,j$\newline（4）$-5j$
\bkjans{（1）\nobreak $3\sqrt3+3j$。（2）\nobreak $-\sqrt2-\sqrt2\,j$。（3）\nobreak $2(\cos120°+j\sin120°)$。（4）\nobreak $5\{\cos(-90°)+j\sin(-90°)\}$。}
\end{jibun}
\end{minipage}\hfill
\begin{minipage}[t]{0.485\linewidth}\vspace{0pt}
\begin{jibun}
$z=4(\cos70°+j\sin70°)$、$w=2(\cos20°+j\sin20°)$ のとき、$zw$ と $\dfrac{z}{w}$ を極形式で求め、$zw$ は $a+bj$ の形にも直しましょう。
\bkjans{$zw=8(\cos90°+j\sin90°)=8j$。$\dfrac{z}{w}=2(\cos50°+j\sin50°)$。}
\end{jibun}
\end{minipage}

\begin{reidai}[かけ算を2通りで計算する]
$(1+j)(\sqrt3+j)$ を、（1）展開して $a+bj$ の形で、（2）極形式で、計算しましょう。また、答えが同じ数か確かめましょう（$\cos75°\approx0.259$、$\sin75°\approx0.966$）。\par
\kangaekata （1）は展開（6-8⑥）。（2）は絶対値と偏角を求めて⑤のきまり。\par
\shiki\hangindent=4.3\zw\hangafter=1 （1）$\sqrt3+j+\sqrt3\,j+j^2=(\sqrt3-1)+(\sqrt3+1)j$\par
\bkvijb $\approx0.732+2.732j$\par
\bkvijq （2）$1+j$ は $\sqrt2$、45°。$\sqrt3+j$ は2、$\tan^{-1}\dfrac{1}{\sqrt3}=30°$\par
\bkvijb 積は、絶対値 $2\sqrt2$、偏角 $45°+30°=75°$\par
\bkvijb 確かめ：$2\sqrt2\cos75°\approx0.732$、$2\sqrt2\sin75°\approx2.732$\par
\kotae\hangindent=4.3\zw\hangafter=1 （1）$(\sqrt3-1)+(\sqrt3+1)j$\newline（2）$2\sqrt2(\cos75°+j\sin75°)$。同じ数
\end{reidai}

\begin{bkfigure}
\begin{tikzpicture}[x=11mm, y=11mm]
  \draw[bkGray, line width=0.6pt, -{Stealth[length=4pt]}] (-1.2,0) -- (2.5,0);
  \draw[bkGray, line width=0.6pt, -{Stealth[length=4pt]}] (0,-0.3) -- (0,3.1);
  \bkvijarrow{bkBlue}{(1,1)}
  \bkvijarrow{bkOrange}{(1.732,1)}
  \bkvijarrow{bkPurple}{(0.732,2.732)}
  \draw[bkOrange, line width=0.6pt] (0.55,0) arc (0:30:0.55);
  \draw[bkBlue, line width=0.6pt] (0.85,0) arc (0:45:0.85);
  \draw[bkPurple, line width=0.6pt] (1.15,0) arc (0:75:1.15);
  \node[font=\sffamily\footnotesize, text=bkOrange!80!black, anchor=west, fill=white, inner sep=1pt] at (1.8,1.0) {$\sqrt3+j$：2、30°};
  \node[font=\sffamily\footnotesize, text=bkBlue!80!black, anchor=east, fill=white, inner sep=1pt] at (-0.15,1.0) {$1+j$：$\sqrt2$、45°};
  \node[font=\sffamily\footnotesize, text=bkPurple, anchor=west, fill=white, inner sep=1pt] at (0.85,2.75) {積：$2\sqrt2$、$45°+30°=75°$};
\end{tikzpicture}
\bkfigcaption{\bkreidairef{5}のかけ算。積の矢印は、長さが $\sqrt2\times2$、角が $45°+30°$。}
\end{bkfigure}

\begin{reidai}[ド・モアブルの定理]
次の数を計算しましょう。\newline（1）$(1+j)^8$\newline（2）$(1+\sqrt3\,j)^6$\newline（3）$(1+j)^4$\newline（4）$(1+\sqrt3\,j)^4$\par
\kangaekata 極形式に直し、絶対値を $n$ 乗、偏角を $n$ 倍する（⑧）。\par
偏角が180°をこえたら、360°をひいて同じ向きの角に直す（4-8の一般角）。\par
\shiki\hangindent=4.3\zw\hangafter=1 （1）$1+j$ は $\sqrt2$、45°。$(\sqrt2)^8=16$、$45°\times8=360°$\par
\bkvijb $16(\cos360°+j\sin360°)=16$\par
\bkvijq （2）$1+\sqrt3\,j$ は2、60°。$2^6=64$、$60°\times6=360°$\par
\bkvijb $64(\cos360°+j\sin360°)=64$\par
\bkvijq （3）$(\sqrt2)^4=4$、$45°\times4=180°$。$4(\cos180°+j\sin180°)=-4$\par
\bkvijq （4）$2^4=16$、$60°\times4=240°$。$240°-360°=-120°$ と同じ向き\par
\bkvijb $16\{\cos(-120°)+j\sin(-120°)\}=16\Bigl(-\dfrac12-\dfrac{\sqrt3}{2}j\Bigr)$\par
\kotae\hangindent=4.3\zw\hangafter=1 （1）16\quad（2）64\quad（3）$-4$\quad（4）$-8-8\sqrt3\,j$
\end{reidai}

\noindent\begin{minipage}[t]{0.485\linewidth}\vspace{0pt}
\begin{jibun}
$(1+\sqrt3\,j)(\sqrt3+j)$ を、展開と極形式の2通りで計算し、同じ数になるか確かめましょう。
\bkjans{展開：$\sqrt3+j+3j+\sqrt3\,j^2=4j$。極形式：$2\times2=4$、$60°+30°=90°$ で $4(\cos90°+j\sin90°)=4j$。同じ。}
\end{jibun}
\end{minipage}\hfill
\begin{minipage}[t]{0.485\linewidth}\vspace{0pt}
\begin{jibun}
次の数を計算しましょう。\par
（1）$(1+j)^6$\newline（2）$(\sqrt3+j)^3$\par
（3）$(1+\sqrt3\,j)^3$
\bkjans{（1）\nobreak $(\sqrt2)^6=8$、$45°\times6=270°$ は $270°-360°=-90°$ の向きなので $-8j$。（2）\nobreak $2^3=8$、90°で $8j$。（3）\nobreak $2^3=8$、180°で $-8$。}
\end{jibun}
\end{minipage}

\begin{reidai}[点を回す]
（1）点 $(4,\ 2)$ を、原点のまわりに左回りに90°回した点を求めましょう。\par
（2）点 $(2,\ 2)$ を、原点のまわりに左回りに30°回した点を求めましょう。\par
\kangaekata $\theta$ 回すには $\cos\theta+j\sin\theta$ をかける（⑨）。90°なら $j$、30°なら $\dfrac{\sqrt3}{2}+\dfrac12j$。\par
\shiki\hangindent=4.3\zw\hangafter=1 （1）$j(4+2j)=4j+2j^2=-2+4j$\par
\bkvijq （2）$(2+2j)\Bigl(\dfrac{\sqrt3}{2}+\dfrac12j\Bigr)=\sqrt3+j+\sqrt3\,j+j^2$\par
\bkvijb $=(\sqrt3-1)+(\sqrt3+1)j$（$j^2=-1$）\par
\bkvijb 確かめ：$|2+2j|=2\sqrt2$、偏角 $45°+30°=75°$。⑤のかけ算の例と同じ点\par
\kotae\hangindent=4.3\zw\hangafter=1 （1）$(-2,\ 4)$\quad（2）$(\sqrt3-1,\ \sqrt3+1)$、約 $(0.73,\ 2.73)$
\end{reidai}

\begin{bkfigure}
\begin{tikzpicture}[x=8mm, y=8mm]
  \draw[bkLine, line width=0.4pt, step=1] (-2.6,-0.6) grid (4.6,4.6);
  \draw[bkGray, line width=0.6pt, -{Stealth[length=4pt]}] (-2.8,0) -- (4.9,0);
  \draw[bkGray, line width=0.6pt, -{Stealth[length=4pt]}] (0,-0.8) -- (0,4.9);
  \bkvijarrow{bkBlue}{(4,2)}
  \bkvijarrow{bkPurple}{(-2,4)}
  \bkvijarrow{bkOrange}{(2,2)}
  \bkvijarrow{bkRed}{(0.732,2.732)}
  \draw[bkGray, line width=0.7pt, dash pattern=on 2pt off 1.5pt, -{Stealth[length=3.5pt]}] (28:3.6) arc (28:115:3.6);
  \draw[bkGray, line width=0.7pt, dash pattern=on 2pt off 1.5pt, -{Stealth[length=3.5pt]}] (47:1.6) arc (47:73:1.6);
  \node[font=\sffamily\footnotesize, anchor=west, fill=white, inner sep=1pt] at (4.3,2) {（1）$(4,\ 2)$};
  \node[font=\sffamily\footnotesize, anchor=south, fill=white, inner sep=1pt] at (-2,4.1) {$(-2,\ 4)$：90°};
  \node[font=\sffamily\footnotesize, anchor=west, fill=white, inner sep=1pt] at (2.2,2.35) {（2）$(2,\ 2)$};
  \node[font=\sffamily\footnotesize, anchor=south west, fill=white, inner sep=1pt] at (0.85,2.8) {30°回した点};
\end{tikzpicture}
\bkfigcaption{\bkreidairef{7}の2つの回転。どちらも、原点からの距離は変わらない。}
\end{bkfigure}

\begin{reidai}[ロボットのうでの先の位置]
4-7の2リンクのうでで、$L_1=\qty{40}{cm}$、$L_2=\qty{30}{cm}$、$\theta_1=30°$、$\theta_2=60°$ のとき、手の位置を、複素数のかけ算で求めましょう（$\sqrt3\approx1.732$）。\par
\kangaekata 手$=u_1(L_1+L_2u_2)$（⑨）。ひじから先を $\theta_2$ 回してから、全体を $\theta_1$ 回す。\par
\shiki\hangindent=2.6\zw\hangafter=1 $u_2=\cos60°+j\sin60°=0.5+0.866j$\par
\bkvija $u_1=\cos30°+j\sin30°=0.866+0.5j$\par
\bkvija $L_1+L_2u_2=40+30\times(0.5+0.866j)=55+25.98j$\par
\bkvija $(55+25.98j)(0.866+0.5j)$\par
\bkvija $=55\times0.866+27.5j+22.50j+12.99j^2$\par
\bkvija $=(47.63-12.99)+(27.5+22.50)j=34.64+50.0j$\par
\kotae\hangindent=2.6\zw\hangafter=1 約 $(34.6,\ 50)$（4-7で求めた位置と同じ）
\end{reidai}

\noindent\begin{minipage}[t]{0.485\linewidth}\vspace{0pt}
\begin{jibun}
（1）点 $(3,\ -1)$ を90°回した点と、（2）点 $(2,\ -2)$ を60°回した点を求めましょう（左回り）。
\bkjans{（1）\nobreak $j(3-j)=1+3j$ で $(1,\ 3)$。（2）\nobreak $(2-2j)\Bigl(\dfrac12+\dfrac{\sqrt3}{2}j\Bigr)=(1+\sqrt3)+(\sqrt3-1)j$ で、約 $(2.73,\ 0.73)$。}
\end{jibun}
\end{minipage}\hfill
\begin{minipage}[t]{0.485\linewidth}\vspace{0pt}
\begin{jibun}
$L_1=\qty{40}{cm}$、$L_2=\qty{30}{cm}$、$\theta_1=60°$、$\theta_2=30°$ のときの手の位置を、複素数のかけ算で求めましょう。
\bkjans{$L_1+L_2u_2=40+30(0.866+0.5j)=65.98+15j$。$u_1=0.5+0.866j$ をかけて $(32.99-12.99)+(57.14+7.5)j=20.0+64.6j$。約 $(20,\ 64.6)$。}
\end{jibun}
\end{minipage}

%% ---------------- まちがえやすいところ・物理工学 ----------------
\bksection{tsumazuki}{まちがえやすいところをチェック}
\noindent\begin{minipage}[t]{0.485\linewidth}\vspace{0pt}
\begin{tsumazuki}{偏角を $\tan^{-1}$ だけで決める}
\bkng{$\arg(-1-j)=45°$}\bkok{$\arg(-1-j)=-135°$}
$\tan^{-1}1=45°$ は右上の向きです。$\tan^{-1}$ は右半分の角しか答えません。$-1-j$ は左下なので $45°-180°=-135°$。先に点をかいて、区画を見ます。
\end{tsumazuki}
\end{minipage}\hfill
\begin{minipage}[t]{0.485\linewidth}\vspace{0pt}
\begin{tsumazuki}{絶対値に $j$ を入れる}
\bkng{$|3+4j|=\sqrt{9+(4j)^2}$}\bkok{$|3+4j|=\sqrt{9+16}=5$}
絶対値は距離なので、実部と虚部（どちらも実数）の2乗をたします。
\end{tsumazuki}
\end{minipage}
\par\vspace{2pt}
\noindent\begin{minipage}[t]{0.485\linewidth}\vspace{0pt}
\begin{tsumazuki}{かけ算で、角をかけてしまう}
\bkng{偏角30°と60°の数の積は、偏角1800°}\bkok{偏角は $30°+60°=90°$}
かけ算では、絶対値はかけ算、偏角は\textbf{たし算}です。
\end{tsumazuki}
\end{minipage}\hfill
\begin{minipage}[t]{0.485\linewidth}\vspace{0pt}
\begin{tsumazuki}{$n$ 乗で、絶対値を忘れる}
\bkng{$(1+j)^8=1$}\bkok{$(1+j)^8=(\sqrt2)^8=16$}
偏角だけ $45°\times8=360°$ にして、絶対値 $\sqrt2$ の8乗を忘れたまちがい。ド・モアブルの定理は絶対値1の数の式です。
\end{tsumazuki}
\end{minipage}

\begin{butsuri}
\bkvijbox{ゲームとロボット：回転はかけ算}
\noindent\begin{minipage}[c]{0.58\linewidth}
\raggedright
ゲームやコンピュータグラフィックスで、キャラクターや絵を回すときは、⑨の「回転をかけ算で表す」考えを使います。平面では、点の1つ1つに $\cos\theta+j\sin\theta$ をかける計算（7-6の回転行列と同じ）です。ロボットのうでのように回転が何回も重なるときも、かけ算を順につなぐだけです。3次元の回転には、複素数をさらに広げた「クォータニオン（四元数）」という数がよく使われます。
\end{minipage}\hfill
\begin{minipage}[c]{0.38\linewidth}
\centering
\begin{tikzpicture}[x=1mm, y=1mm]
  \draw[bkGray, line width=0.5pt] (-4,0) -- (32,0) (0,-4) -- (0,30);
  \foreach \a/\c in {0/bkBlue, 30/bkOrange, 60/bkPurple} {
    \begin{scope}[rotate=\a]
      \filldraw[fill=\c!25, draw=\c, line width=0.9pt] (14,-3.5) -- (26,0) -- (14,3.5) -- (17,0) -- cycle;
    \end{scope}}
  \draw[bkGray, line width=0.7pt, dash pattern=on 2pt off 1.5pt, -{Stealth[length=3.5pt]}] (5:10) arc (5:62:10);
  \node[font=\sffamily\footnotesize, anchor=north] at (15,-4.5) {30°ずつ回す};
\end{tikzpicture}
\end{minipage}
\end{butsuri}

\begin{butsuri}
\bkvijbox{交流：回る矢印の影（6-10の予告）}
\noindent\begin{minipage}[c]{0.58\linewidth}
\raggedright
家のコンセントの交流の電圧 $V_0\cos\omega t$ は、絶対値 $V_0$、偏角 $\omega t$ の複素数、つまり角速度 $\omega$ で回る矢印の、実軸への影（実部）と見ることができます（4-8の「円運動の影」）。電気の本では、この回る矢印を\textbf{フェーザ}とよびます。ずれのある2つの交流をたすことは、矢印をつなぐ（②）ことになり、計算がとても楽になります。くわしくは6-10で学びます。
\end{minipage}\hfill
\begin{minipage}[c]{0.38\linewidth}
\centering
\begin{tikzpicture}[x=1mm, y=1mm]
  \draw[bkGray, line width=0.5pt, -{Stealth[length=3pt]}] (-17,0) -- (19,0) node[anchor=north east, font=\sffamily\footnotesize, inner sep=1pt, yshift=-1mm] {実軸};
  \draw[bkGray, line width=0.5pt] (0,-16) -- (0,16);
  \draw[bkInk!35, line width=0.6pt] (0,0) circle (13);
  \draw[bkBlue, line width=1.4pt, -{Stealth[length=4pt]}] (0,0) -- (50:13);
  \draw[bkGray, line width=0.6pt, dash pattern=on 2pt off 1.5pt] (50:13) -- ({13*cos(50)},0);
  \draw[bkOrange, line width=2pt] (0,0) -- ({13*cos(50)},0);
  \draw[bkInk, line width=0.5pt] (4,0) arc (0:50:4);
  \node[font=\sffamily\footnotesize] at (22:6.2) {$\omega t$};
  \node[font=\sffamily\footnotesize, text=bkBlue!80!black, anchor=south west, inner sep=1pt] at (50:13.4) {$V_0$};
  \node[font=\sffamily\footnotesize, text=bkOrange!80!black, anchor=north, inner sep=1pt, align=center, fill=bkPurpleBg] at (5,-2.2) {影\\$V_0\cos\omega t$};
  \draw[bkInk!50, line width=0.5pt, -{Stealth[length=2.5pt]}] (60:15.5) arc (60:95:15.5);
\end{tikzpicture}
\end{minipage}
\end{butsuri}

\begin{butsuri}
\bkvijbox{三相交流とモーター}
\noindent\begin{minipage}[c]{0.58\linewidth}
\raggedright
発電所から工場へ送る電気は、3本の電線の電圧が120°ずつずれた「三相交流」です（4-9）。3つの電圧を回る矢印で表すと、たがいに120°ずつ回った3本の矢印で、つなぐともとにもどるので、和はいつも0です（1の3乗根。もっと知りたい）。工場のモーターでは、120°ずつずらして置いた3つのコイルにこの電流を流して、ぐるぐる回る磁石のはたらき（回転磁界）を作り、モーターを回しています（物理シリーズ『電磁気学』）。
\end{minipage}\hfill
\begin{minipage}[c]{0.38\linewidth}
\centering
\begin{tikzpicture}[x=1mm, y=1mm]
  \draw[bkGray, line width=0.5pt] (-13,0) -- (13,0) (0,-13) -- (0,15);
  \draw[bkRed, line width=1.4pt, -{Stealth[length=4pt]}] (0,0) -- (90:12);
  \draw[bkBlue, line width=1.4pt, -{Stealth[length=4pt]}] (0,0) -- (210:12);
  \draw[bkGreen!70!black, line width=1.4pt, -{Stealth[length=4pt]}] (0,0) -- (330:12);
  \node[font=\sffamily\footnotesize, anchor=north, inner sep=1pt] at (0,-14) {120°ずつずれた3本};
  \begin{scope}[xshift=26mm, yshift=-4mm]
    \draw[bkRed, line width=1.4pt, -{Stealth[length=4pt]}] (0,0) -- (90:12);
    \draw[bkBlue, line width=1.4pt, -{Stealth[length=4pt]}] (90:12) -- ($(90:12)+(210:12)$);
    \draw[bkGreen!70!black, line width=1.4pt, -{Stealth[length=4pt]}] ($(90:12)+(210:12)$) -- (0,0);
    \node[font=\sffamily\footnotesize, anchor=north, inner sep=1pt, align=center] at (-3,-2) {つなぐと\\もとにもどる};
  \end{scope}
\end{tikzpicture}
\end{minipage}
\end{butsuri}

\begin{column}{「ないはずの数」が、平面の点になるまで}
虚数は、はじめ「計算のとちゅうに出てくる、ふしぎな記号」でした（6-8のコラム）。それを平面の点として目に見える形にしたのが、ノルウェー生まれで、デンマークで測量の仕事をしていたウェッセルです（1797年に発表、1799年に論文）。デンマーク語で書かれていたので、長いあいだ知られませんでした。1806年には、パリのアルガンが同じ考えを本にしました。大数学者ガウスが1831年に「複素数」という名前とともにこの考えを示して、広く知られるようになりました。
\par\vspace{3pt}
アイルランドのハミルトンは、1830年代に複素数を「2つの実数の組」として決めなおし、さらに3次元の回転を表せる数（クォータニオン）を考え出しました（1843年）。いまでは、ドローンや宇宙船の向きの計算、ゲームの3次元の画面に使われています。
\end{column}

%% ---------------- 章末問題 ----------------
\bkendproblems
\bkproblemgroup{基本}{4}

\bkproblem 次の複素数の絶対値と偏角を求めましょう\mbox{（\bkreidairef{1}）}。\newline\bkvijit{（1）$-3+3j$}\bkvijit{（2）$1-\sqrt3\,j$}\bkvijit{（3）$-4$}
\bkpans{（1）\nobreak $3\sqrt2$、$-45°+180°=135°$。（2）\nobreak 2、$-60°$。（3）\nobreak 4、180°。}

\bkproblem （1）$4(\cos120°+j\sin120°)$ を $a+bj$ の形に、（2）$1-j$ を極形式に直しましょう\mbox{（\bkreidairef{3}）}。
\bkpans{（1）\nobreak $-2+2\sqrt3\,j$。（2）\nobreak $\sqrt2\{\cos(-45°)+j\sin(-45°)\}$。}

\bkproblem $z_1=-1+2j$、$z_2=5+10j$ のとき、$z_1+z_2$、$z_2-z_1$ と、2つの点の距離を求めましょう\mbox{（\bkreidairef{2}）}。
\bkpans{$z_1+z_2=4+12j$、$z_2-z_1=6+8j$、距離 $\sqrt{36+64}=10$。}

\bkproblem $z=3(\cos40°+j\sin40°)$、$w=2(\cos50°+j\sin50°)$ のとき、$zw$ と $\dfrac{z}{w}$ を極形式で求め、$zw$ は $a+bj$ の形にも直しましょう\mbox{（\bkreidairef{4}）}。
\bkpans{$zw=6(\cos90°+j\sin90°)=6j$。$\dfrac{z}{w}=\dfrac32\{\cos(-10°)+j\sin(-10°)\}$。}

\bkproblem $(1+\sqrt3\,j)(1+j)$ を、展開と極形式の2通りで求めましょう\mbox{（\bkreidairef{5}）}。
\bkpans{展開：$1+j+\sqrt3\,j+\sqrt3\,j^2=(1-\sqrt3)+(1+\sqrt3)j$。極形式：$2\times\sqrt2=2\sqrt2$、$60°+45°=105°$ で $2\sqrt2(\cos105°+j\sin105°)$。}

\bkproblem 次の数を計算しましょう\mbox{（\bkreidairef{6}）}。\newline\bkvijit{（1）$(1+j)^{10}$}\bkvijit{（2）$(-1+\sqrt3\,j)^3$}
\bkpans{（1）\nobreak $(\sqrt2)^{10}=32$、$45°\times10=450°$、$450°-360°=90°$ なので $32j$。（2）\nobreak $2^3=8$、$120°\times3=360°$ で8。}

\bkproblem 点 $(5,\ 1)$ を、原点のまわりに左回りに90°回した点を求めましょう\mbox{（\bkreidairef{7}）}。
\bkpans{$j(5+j)=5j+j^2=-1+5j$ で $(-1,\ 5)$。}

\bkproblem $-1+j$ を極形式に直しましょう\mbox{（\bkreidairef{3}）}。
\bkpans{$r=\sqrt2$、左上で $\tan^{-1}(-1)=-45°$、$-45°+180°=135°$。$\sqrt2(\cos135°+j\sin135°)$。}

\bkproblemgroup{チャレンジ}{5}

\bkproblem 点 $(-2,\ 4)$ を、原点のまわりに左回りに（1）90°、（2）45° 回した点を求めましょう\mbox{（\bkreidairef{7}）}。
\bkpans{（1）\nobreak $j(-2+4j)=-4-2j$ で $(-4,\ -2)$。（2）\nobreak $(-2+4j)\cdot\dfrac{\sqrt2}{2}(1+j)=\dfrac{\sqrt2}{2}(-6+2j)$ で $(-3\sqrt2,\ \sqrt2)$、約 $(-4.24,\ 1.41)$。}

\bkproblem 2リンクのうでで、$L_1=\qty{40}{cm}$、$L_2=\qty{30}{cm}$、$\theta_1=90°$、$\theta_2=-90°$ のときの手の位置を、複素数のかけ算で求めましょう\mbox{（\bkreidairef{8}）}。
\bkpans{$u_1=j$、$u_2=-j$。$j(40+30\times(-j))=j(40-30j)=30+40j$。$(30,\ 40)$。}

\bkproblem $|z|=2$、$\arg z=30°$ の複素数 $z$ について、$z^2$、$z^3$、$z^6$ を $a+bj$ の形で求めましょう\mbox{（\bkreidairef{6}）}。
\bkpans{$z^2=4(\cos60°+j\sin60°)=2+2\sqrt3\,j$。$z^3=8(\cos90°+j\sin90°)=8j$。$z^6=64(\cos180°+j\sin180°)=-64$。}

\bkproblem $\dfrac{z}{w}$ を求めましょう\mbox{（\bkreidairef{4}）}。\newline（1）$z=4(\cos100°+j\sin100°)$、$w=2(\cos10°+j\sin10°)$ のとき、極形式と $a+bj$ の形で。\newline（2）$z=1+\sqrt3\,j$、$w=1+j$ のとき、極形式で。また、分母と共役な数をかける方法（6-8）で $a+bj$ の形にも直し、絶対値が同じになるか確かめましょう。
\bkpans{（1）\nobreak 絶対値 $\dfrac42=2$、偏角 $100°-10°=90°$ で $2(\cos90°+j\sin90°)=2j$。（2）\nobreak $|z|=2$、60°、$|w|=\sqrt2$、45° なので $\sqrt2(\cos15°+j\sin15°)$。$a+bj$ の形は $\dfrac{(1+\sqrt3)+(\sqrt3-1)j}{2}$、約 $1.366+0.366j$。絶対値は $\sqrt{1.866+0.134}=\sqrt2$ で同じ。}

\begin{furikaeri}
\begin{checklist}
\item 複素数を平面の点で表し、絶対値と偏角を求められる。\bkfurref{bk:6-9:kyoku}
\item かけ算が「大きさはかけ算、角はたし算」になることを説明できる。\bkfurref{bk:6-9:kakeru}
\item かけ算で、点を回したり、うでの先の位置を求めたりできる。\bkfurref{bk:6-9:kaiten}
\end{checklist}
\end{furikaeri}

\bkfukidashi{mina}{かけ算が、回転と拡大になるなんて！　次の6-10では、$e$ と回転がつながる「オイラーの公式」に進もう。}
